\documentclass[a4paper]{article}

% ── Paquetes ──────────────────────────────────────────────────────────────────
\usepackage[margin=0.6cm]{geometry}
\usepackage[T1]{fontenc}
\usepackage[utf8]{inputenc}
\usepackage[spanish]{babel}
\usepackage{tikz}
\usepackage{xcolor}
\usepackage{fontawesome5}
\usepackage{lmodern}
\usepackage{microtype}
\usepackage{mdframed}
\usepackage{graphicx}
\usepackage{array}
\usepackage{booktabs}
\usepackage{multirow}
\usetikzlibrary{shapes.geometric, arrows.meta, positioning, shadows, backgrounds, calc}

% ── Paleta de colores ──────────────────────────────────────────────────────────
\definecolor{netflix}{HTML}{E50914}
\definecolor{darkbg}{HTML}{141414}
\definecolor{cardbg}{HTML}{1F1F1F}
\definecolor{accent}{HTML}{E50914}
\definecolor{gold}{HTML}{F5A623}
\definecolor{silver}{HTML}{AAAAAA}
\definecolor{lightgray}{HTML}{2A2A2A}
\definecolor{textwhite}{HTML}{FFFFFF}
\definecolor{textgray}{HTML}{CCCCCC}
\definecolor{highlight}{HTML}{FF6B6B}

\pagecolor{darkbg}
\color{textwhite}

% ── Sin número de página ──────────────────────────────────────────────────────
\pagestyle{empty}

\begin{document}

% ════════════════════════════════════════════════════════════════════════════════
%  CABECERA PRINCIPAL
% ════════════════════════════════════════════════════════════════════════════════
\begin{tikzpicture}[remember picture, overlay]
  % Barra superior roja
  \fill[netflix] (current page.north west) rectangle +(\paperwidth, -2.8cm);

  % Etiqueta 5V
  \node[anchor=west, text=textwhite, font=\fontsize{9}{9}\selectfont\bfseries,
        inner sep=0pt] at ([xshift=1cm, yshift=-0.55cm]current page.north west)
    {\textls[200]{LAS 5 V DEL BIG DATA}};

  % Título grande
  \node[anchor=west, text=textwhite, font=\fontsize{34}{34}\selectfont\bfseries,
        inner sep=0pt] at ([xshift=1cm, yshift=-1.45cm]current page.north west)
    {VALOR};

  % Subtítulo
  \node[anchor=west, text=textwhite!80, font=\fontsize{11}{11}\selectfont,
        inner sep=0pt] at ([xshift=1cm, yshift=-2.25cm]current page.north west)
    {Caso de éxito: \textbf{Netflix} y la producción de \textit{House of Cards}};

  % Ícono lado derecho
  \node[anchor=east, text=textwhite, font=\fontsize{42}{42}\selectfont,
        inner sep=0pt] at ([xshift=-1cm, yshift=-1.4cm]current page.north east)
    {\faChartLine};
\end{tikzpicture}

\vspace{2.9cm}

% ════════════════════════════════════════════════════════════════════════════════
%  FILA 1 — Definición + Contexto
% ════════════════════════════════════════════════════════════════════════════════
\noindent
\begin{tikzpicture}
% ── Card: Qué es el Valor ──────────────────────────────────────────────────────
\node[anchor=north west,
      fill=cardbg,
      rounded corners=6pt,
      minimum width=8.6cm,
      minimum height=4.6cm,
      text width=8cm,
      inner sep=12pt] (defcard) at (0,0) {
  {\color{netflix}\faDatabase\quad\fontsize{10}{10}\selectfont\bfseries¿QUÉ ES EL VALOR EN BIG DATA?}\\[6pt]
  {\color{textgray}\fontsize{9.5}{13}\selectfont
  El \textbf{\color{textwhite}Valor} es la capacidad de transformar
  grandes volúmenes de datos crudos en \textbf{\color{gold}información
  accionable} que genere un impacto económico, social o estratégico
  concreto y medible.\\[6pt]
  No basta con tener datos: deben convertirse en
  \textbf{\color{textwhite}decisiones}.}
};

% ── Card: Contexto ────────────────────────────────────────────────────────────
\node[anchor=north west,
      fill=cardbg,
      rounded corners=6pt,
      minimum width=9.6cm,
      minimum height=4.6cm,
      text width=9cm,
      inner sep=12pt] (ctxcard) at (9cm,0) {
  {\color{netflix}\faFilm\quad\fontsize{10}{10}\selectfont\bfseries CONTEXTO}\\[6pt]
  {\color{textgray}\fontsize{9.5}{13}\selectfont
  En \textbf{\color{textwhite}2011}, Netflix contaba con
  \textbf{\color{gold}+30 millones} de suscriptores y acumulaba datos de
  comportamiento masivos: géneros preferidos, actores y directores con mayor
  retención, episodios donde los usuarios abandonaban y patrones de
  consumo compulsivo (\textit{binge-watching}).\\[6pt]
  Estos datos eran la materia prima de una decisión de
  \textbf{\color{netflix}\$100.000.000}.}
};
\end{tikzpicture}

\vspace{0.35cm}

% ════════════════════════════════════════════════════════════════════════════════
%  FILA 2 — Cómo se tomó la decisión (pipeline de datos)
% ════════════════════════════════════════════════════════════════════════════════
\noindent
\begin{tikzpicture}
\node[anchor=north west,
      fill=cardbg,
      rounded corners=6pt,
      minimum width=18.6cm,
      text width=18cm,
      inner sep=14pt] (pipeline) at (0,0) {
  {\color{netflix}\faCodeBranch\quad\fontsize{10}{10}\selectfont\bfseries CRUCE DE DATOS QUE RESPALDÓ LA INVERSIÓN}\\[10pt]
  % Diagrama de flujo con TikZ interno
  \begin{tikzpicture}[
    box/.style={fill=lightgray, rounded corners=5pt, minimum width=4.2cm,
                minimum height=1.5cm, text width=3.8cm, align=center,
                font=\fontsize{8.5}{11}\selectfont, text=textgray, inner sep=6pt},
    arrow/.style={-Stealth, thick, color=netflix},
    node distance=0.6cm
  ]
    \node[box] (v1) {
      {\color{gold}\faUsers}\\[3pt]
      Usuarios que vieron\\
      \textbf{\color{textwhite}House of Cards (1990)}\\
      y la valoraron positivamente
    };
    \node[box, right=of v1] (v2) {
      {\color{gold}\faUserTie}\\[3pt]
      De ese grupo, los que\\
      consumían films de\\
      \textbf{\color{textwhite}Kevin Spacey}
    };
    \node[box, right=of v2] (v3) {
      {\color{gold}\faVideo}\\[3pt]
      De ese grupo, los que\\
      valoraban obras de\\
      \textbf{\color{textwhite}David Fincher}
    };

    \node[fill=netflix, rounded corners=5pt, minimum width=4.2cm,
          minimum height=1.5cm, text width=3.8cm, align=center,
          font=\fontsize{8.5}{11}\selectfont, text=textwhite, inner sep=6pt,
          right=of v3] (resultado) {
      {\faChartPie}\\[3pt]
      \textbf{Solapamiento}\\
      estadísticamente\\
      significativo
    };

    \draw[arrow] (v1) -- (v2);
    \draw[arrow] (v2) -- (v3);
    \draw[arrow] (v3) -- (resultado);
  \end{tikzpicture}
  \\[8pt]
  {\color{textgray}\fontsize{9}{12}\selectfont
  El análisis cruzado identificó una \textbf{\color{textwhite}masa crítica de suscriptores activos}
  con alta probabilidad de consumir el contenido \textit{antes} de producirlo, convirtiendo
  la inversión en una apuesta con respaldo cuantitativo.}
};
\end{tikzpicture}

\vspace{0.35cm}

% ════════════════════════════════════════════════════════════════════════════════
%  FILA 3 — KPIs + Decisión + Expansión
% ════════════════════════════════════════════════════════════════════════════════
\noindent
\begin{tikzpicture}

% ── KPIs ──────────────────────────────────────────────────────────────────────
\node[anchor=north west,
      fill=cardbg,
      rounded corners=6pt,
      minimum width=5.8cm,
      text width=5.2cm,
      inner sep=13pt] (kpi) at (0,0) {
  {\color{netflix}\faTrophy\quad\fontsize{10}{10}\selectfont\bfseries RESULTADOS}\\[8pt]
  % KPI 1
  \begin{tikzpicture}
    \fill[netflix] (0,0) rectangle (0.18cm, 0.7cm);
    \node[anchor=west, text=textgray, font=\fontsize{8.5}{11}\selectfont,
          text width=4.2cm] at (0.28cm, 0.35cm)
      {\textbf{\color{gold}9 nominaciones} al Emmy en la primera temporada.};
  \end{tikzpicture}\\[5pt]
  \begin{tikzpicture}
    \fill[netflix] (0,0) rectangle (0.18cm, 0.7cm);
    \node[anchor=west, text=textgray, font=\fontsize{8.5}{11}\selectfont,
          text width=4.2cm] at (0.28cm, 0.35cm)
      {\textbf{\color{gold}Reducción del churn} (cancelaciones) en el período de lanzamiento.};
  \end{tikzpicture}\\[5pt]
  \begin{tikzpicture}
    \fill[netflix] (0,0) rectangle (0.18cm, 0.7cm);
    \node[anchor=west, text=textgray, font=\fontsize{8.5}{11}\selectfont,
          text width=4.2cm] at (0.28cm, 0.35cm)
      {\textbf{\color{gold}Nuevas suscripciones} en mercados de alta exposición.};
  \end{tikzpicture}\\[5pt]
  \begin{tikzpicture}
    \fill[netflix] (0,0) rectangle (0.18cm, 0.7cm);
    \node[anchor=west, text=textgray, font=\fontsize{8.5}{11}\selectfont,
          text width=4.2cm] at (0.28cm, 0.35cm)
      {\textbf{\color{gold}80\% del catálogo} original Netflix se decide hoy con este modelo.};
  \end{tikzpicture}
};

% ── Decisión económica ────────────────────────────────────────────────────────
\node[anchor=north west,
      fill=netflix,
      rounded corners=6pt,
      minimum width=6.6cm,
      text width=6cm,
      inner sep=14pt] (decision) at (6.3cm,0) {
  {\color{textwhite}\faDollarSign\quad\fontsize{10}{10}\selectfont\bfseries DECISIÓN ECONÓMICA}\\[8pt]
  {\fontsize{28}{28}\selectfont\bfseries\color{textwhite}\$100M}\\
  {\fontsize{8.5}{12}\selectfont\color{textwhite!85}
  de inversión en producción propia \textbf{sin piloto} de prueba ni
  evaluación de \textit{focus groups}, respaldada \textbf{100\% por datos}
  de comportamiento interno.\\[6pt]
  Esto eliminó el ciclo tradicional de \textbf{meses} de evaluación
  y lo redujo a \textbf{semanas}.}
};

% ── Valor vs Método Tradicional ───────────────────────────────────────────────
\node[anchor=north west,
      fill=cardbg,
      rounded corners=6pt,
      minimum width=5.9cm,
      text width=5.3cm,
      inner sep=13pt] (compare) at (13.4cm,0) {
  {\color{netflix}\faExchangeAlt\quad\fontsize{10}{10}\selectfont\bfseries MODELO TRADICIONAL vs. BIG DATA}\\[8pt]
  {\color{textgray}\fontsize{8.5}{13}\selectfont
  \textbf{\color{silver}Antes (sin datos):}\\
  Piloto → Focus group → Aprobación editorial → Producción.
  \textbf{Alta incertidumbre.}\\[6pt]
  \textbf{\color{gold}Con Big Data:}\\
  Análisis de comportamiento → Audiencia pre-identificada → Decisión fundamentada.
  \textbf{Riesgo cuantificado.}}
};

\end{tikzpicture}

\vspace{0.35cm}

% ════════════════════════════════════════════════════════════════════════════════
%  FILA 4 — Conclusión
% ════════════════════════════════════════════════════════════════════════════════
\noindent
\begin{tikzpicture}
\node[anchor=north west,
      fill=lightgray,
      rounded corners=6pt,
      minimum width=18.6cm,
      text width=18cm,
      inner sep=14pt] (conclusion) at (0,0) {
  {\color{netflix}\faLightbulb\quad\fontsize{10}{10}\selectfont\bfseries CONCLUSIÓN — LA V DE VALOR}\\[6pt]
  {\color=textgray\fontsize{9.5}{14}\selectfont
  El \textbf{\color{textwhite}Valor} en Big Data no reside en el volumen de información almacenada,
  sino en la capacidad de extraer de ella una \textbf{\color{gold}decisión económica concreta y medible}.
  Netflix transformó datos de comportamiento —\textit{clicks}, historial, ratings— en una inversión de
  \textbf{\color{netflix}\$100 millones} con respaldo estadístico, desplazando la intuición editorial
  por \textbf{\color{textwhite}inteligencia basada en evidencia}. Ese es el Valor real del Big Data.}
};
\end{tikzpicture}

\vspace{0.3cm}

% ════════════════════════════════════════════════════════════════════════════════
%  PIE DE PÁGINA
% ════════════════════════════════════════════════════════════════════════════════
\noindent
\begin{tikzpicture}
  \draw[silver!30, line width=0.4pt] (0,0) -- (18.6cm, 0);
  \node[anchor=west, font=\fontsize{7.5}{7.5}\selectfont, text=silver]
    at (0, -0.3cm) {Fuente: Netflix Tech Blog · Harvard Business Review — Caso Netflix Big Data Strategy (2013-2014)};
  \node[anchor=east, font=\fontsize{7.5}{7.5}\selectfont, text=silver]
    at (18.6cm, -0.3cm) {Big Data · 5V · Valor};
\end{tikzpicture}

\end{document}
